% !TEX encoding = UTF-8 Unicode
% !TEX TS-program = pdflatex

% toptesi document class
%\documentclass[%
%    a4paper, % not needed, by default it is a4paper, or also b5paper can be used
%    corpo=12pt, % dimension of basic font
%    % oneside is generally the way to go
%    oneside, % two side optimizes for two-face printing, having chapters open on the right (aka odd numbers), if you don't want blank pages put oneside here
%    stile=standard,
%    %evenboxes, % not needed, to put supervisors and candidate at the same level
%    tipotesi=magistrale,
%    numerazioneromana, % roman numbering for appendixes and preambles, up to Table of Contents
%    openright, % to force opening on the right for double-sided printing
%    cucitura=7mm, % for printing, 7mm should be enough
    %dvipsnames, % for compatibility with xcolor, it does not work
%]{toptesi}

\documentclass[11pt]{report}


\setcounter{secnumdepth}{4}
\setcounter{tocdepth}{4}

\usepackage{calc} % needed for complex computationsin latex language
% \usepackage[a4paper]{geometry} % needed to change margins


% configuration for glossaries
% convert and load converted glossaries in .tex ,format from .bib
\setabbreviationstyle{long-short-desc} % style before loading resources
% this command sets the style to title for long names of acronyms only in the glossary description, leading to capitalized first-letter for all words
% \glssetcategoryattribute{\glsxtrabbrvtype}{glossname}{capitalisewords} % doesn't work
% resources to load if using a bib file with bib2gls
%\GlsXtrLoadResources[%
% src={glossaries}, % name of the file without extension
% selection=all, % select all the entries
%]
% not needed
%\newglossary*{abbreviation}{Acronyms} % to change the name of this glossary for acronyms

%\renewcommand{\glsclearpage}{\paginavuota} % to allow glossaries to clear pages, done manually is better

% setup for hyperref
\hypersetup{%
    pdfpagemode={UseOutlines},
    bookmarksopen,
    pdfstartview={FitH},
    colorlinks,
    linkcolor={black}, % it is suggested to keep them black, since when printing it it costs per page, and if they have color it's twice the price per page
    citecolor={black},
    urlcolor={black}
  }
%

% setup for svg
\svgsetup{%
    inkscapeformat=pdf, % to force usage of PDF
    inkscapelatex=false, % to disable latex rendering of text, produces errors
}

% setup for siunitx, it does not work in the summary
\sisetup{%
    detect-all, % to use the same font as for writing when using \num
    mode=text, % to allow it to work also in math mode
    group-separator = {,}, % separator for number grouping
    group-minimum-digits = 3, % minimum number of digits a number must have to be grouped in 3-digit groups
}

% listings colours
\definecolor{rulecolor}{rgb}{0,0,0}
\definecolor{commentcolor}{rgb}{0,0.6,0}
\definecolor{linenumbercolor}{rgb}{0.5,0.5,0.5}
\definecolor{keywordcolor}{rgb}{0,0,0.95}
\definecolor{backcolor}{rgb}{1,1,1}%{0.95,0.95,0.92}
\definecolor{stringcolor}{rgb}{0.58,0,0.82}

% setup for lstlisting
\lstset{ %
	backgroundcolor=\color{backcolor},   % choose the background color; you must add \usepackage{color} or \usepackage{xcolor}; should come as last argument
	basicstyle=\footnotesize,        % the size of the fonts that are used for the code
	breakatwhitespace=false,         % sets if automatic breaks should only happen at whitespace
	breaklines=true,                 % sets automatic line breaking
	captionpos=t,                    % sets the caption-position to bottom
	commentstyle=\color{commentcolor},    % comment style
	extendedchars=true,              % lets you use non-ASCII characters; for 8-bits encodings only, does not work with UTF-8
	frame=single,	                   % adds a frame around the code
	keepspaces=true,                 % keeps spaces in text, useful for keeping indentation of code (possibly needs columns=flexible)
	keywordstyle=\color{keywordcolor},       % keyword style
	%language=VHDL,                 % the language of the code
	numbers=left,                    % where to put the line-numbers; possible values are (none, left, right)
	numbersep=5pt,                   % how far the line-numbers are from the code
	numberstyle=\tiny\color{linenumbercolor}, % the style that is used for the line-numbers
	rulecolor=\color{rulecolor},         % if not set, the frame-color may be changed on line-breaks within not-black text (e.g. comments (green here))
	showspaces=false,                % show spaces everywhere adding particular underscores; it overrides 'showstringspaces'
	showstringspaces=false,          % underline spaces within strings only
	showtabs=false,                  % show tabs within strings adding particular underscores
	stepnumber=1,                    % the step between two line-numbers. If it's 1, each line will be numbered
	stringstyle=\color{stringcolor},     % string literal style
	tabsize=4,	                   % sets default tabsize to 2 spaces
	title=\lstname,                   % show the filename of files included with \lstinputlisting; also try caption instead of title
	inputencoding=utf8,
	literate=
	{á}{{\'a}}1 {é}{{\'e}}1 {í}{{\'i}}1 {ó}{{\'o}}1 {ú}{{\'u}}1
	{Á}{{\'A}}1 {É}{{\'E}}1 {Í}{{\'I}}1 {Ó}{{\'O}}1 {Ú}{{\'U}}1
	{à}{{\`a}}1 {è}{{\`e}}1 {ì}{{\`i}}1 {ò}{{\`o}}1 {ù}{{\`u}}1
	{À}{{\`A}}1 {È}{{\'E}}1 {Ì}{{\`I}}1 {Ò}{{\`O}}1 {Ù}{{\`U}}1
	{ä}{{\"a}}1 {ë}{{\"e}}1 {ï}{{\"i}}1 {ö}{{\"o}}1 {ü}{{\"u}}1
	{Ä}{{\"A}}1 {Ë}{{\"E}}1 {Ï}{{\"I}}1 {Ö}{{\"O}}1 {Ü}{{\"U}}1
	{â}{{\^a}}1 {ê}{{\^e}}1 {î}{{\^i}}1 {ô}{{\^o}}1 {û}{{\^u}}1
	{Â}{{\^A}}1 {Ê}{{\^E}}1 {Î}{{\^I}}1 {Ô}{{\^O}}1 {Û}{{\^U}}1
	{œ}{{\oe}}1 {Œ}{{\OE}}1 {æ}{{\ae}}1 {Æ}{{\AE}}1 {ß}{{\ss}}1
	{ű}{{\H{u}}}1 {Ű}{{\H{U}}}1 {ő}{{\H{o}}}1 {Ő}{{\H{O}}}1
	{ç}{{\c c}}1 {Ç}{{\c C}}1 {ø}{{\o}}1 {å}{{\r a}}1 {Å}{{\r A}}1
	{€}{{\euro}}1 {£}{{\pounds}}1 {«}{{\guillemotleft}}1
	{»}{{\guillemotright}}1 {ñ}{{\~n}}1 {Ñ}{{\~N}}1 {¿}{{?`}}1
}


% biblatex setup
% generally 9000 is ok, if higher than 10000 it's bad
% If you want to break on URL numbers
\setcounter{biburlnumpenalty}{9000}
% If you want to break on URL lower case letters
\setcounter{biburllcpenalty}{9000}
% If you want to break on URL UPPER CASE letters
\setcounter{biburlucpenalty}{9000}

\definecolor{mGreen}{rgb}{0,0.6,0}
\definecolor{mGray}{rgb}{0.5,0.5,0.5}
\definecolor{mPurple}{rgb}{0.58,0,0.82}
\definecolor{backgroundColour}{rgb}{0.95,0.95,0.92}

\lstdefinestyle{CStyle}{
    backgroundcolor=\color{backgroundColour},   
    commentstyle=\color{mGreen},
    keywordstyle=\color{magenta},
    numberstyle=\tiny\color{mGray},
    stringstyle=\color{mPurple},
    basicstyle=\footnotesize\ttfamily,
    breakatwhitespace=false,         
    breaklines=true,                 
    captionpos=b,                    
    keepspaces=true,                 
    numbers=left,                    
    numbersep=5pt,                  
    showspaces=false,                
    showstringspaces=false,
    showtabs=false,                  
    tabsize=2,
    language=C,
    lineskip=-0.1ex
}


% how to change Contents to Table of Contents
\addto\captionsenglish{% Replace "english" with the language you use
  \renewcommand{\contentsname}%
    {Table of Contents}%
}

% to change the name of Abbreviations to Acronyms
% not needed if use use entry types and define those
% \renewcommand{\abbreviationsname}{Acronyms}

% to allow line comments in algorithms
\algnewcommand{\LineComment}[1]{\State \(\triangleright\) #1}

% to declare abs and norm
\DeclarePairedDelimiter\abs{\lvert}{\rvert}%
\DeclarePairedDelimiter\norm{\lVert}{\rVert}%

% Swap the definition of \abs* and \norm*, so that \abs
% and \norm resizes the size of the brackets, and the 
% starred version does not.
\makeatletter
\let\oldabs\abs
\def\abs{\@ifstar{\oldabs}{\oldabs*}}
%
\let\oldnorm\norm
\def\norm{\@ifstar{\oldnorm}{\oldnorm*}}
\makeatother


\lstdefinestyle{CStyle}{
    backgroundcolor=\color{backgroundColour},   
    commentstyle=\color{mGreen},
    keywordstyle=\color{magenta},
    numberstyle=\tiny\color{mGray},
    stringstyle=\color{mPurple},
    basicstyle=\scriptsize\ttfamily,
    breakatwhitespace=false,         
    breaklines=true,                 
    captionpos=b,                    
    keepspaces=true,                 
    numbers=left,                    
    numbersep=5pt,                  
    showspaces=false,                
    showstringspaces=false,
    showtabs=false,                  
    tabsize=2,
    language=C,
    lineskip=-0.1ex
}

\newcommand{\Dedications}{%
    \chapter*{Dedications}
    \addcontentsline{toc}{chapter}{Dedications} % Aggiunge 'Dedications' all'indice
}

\newcommand{\paginavuota}{
    \newpage
    \thispagestyle{empty}
}

% Formato per subsubsection
\titleformat{\subsubsection}[block]
  {\normalfont\normalsize\bfseries} % Stile del titolo
  {\thesubsubsection} % Aggiungi il numero di subsubsection
  {1em} % Spaziatura tra numero e titolo
  {} % Comando dopo il titolo

\renewcommand\thesubsubsection{\thesection.\arabic{subsection}.\arabic{subsubsection}} % Formato della numerazione per subsubsection

% Formato per paragraph
\titleformat{\paragraph}[runin]
  {\normalfont\normalsize\bfseries} % Stile del titolo
  {\theparagraph} % Aggiungi il numero di paragraph
  {1em} % Spaziatura tra numero e titolo
  {} % Comando dopo il titolo

\renewcommand\theparagraph{\thesubsubsection.\arabic{paragraph}} % Formato della numerazione per paragraph

% Definisci un nuovo comando per la subsubsubsection
% Definisci un nuovo comando per la subsubsubsection
\newcommand{\subsubsubsection}[1]{%
  \paragraph{#1}\mbox{}\\ % Usa la formattazione del paragraph e va a capo
  \addcontentsline{toc}{paragraph}{#1} % Aggiungi al sommario
}



% change this configuration with your info
% if you need fewer or more supervisors you have to change \relatore command by adding or removing lines in the table in toptesi_config
\newcommand{\thesistitle}{Thesis TITLE}
\newcommand{\thesisuniversitylogo}{images/polito_logo_2021_blu} % choose your logo
\newcommand{\thesiscandidatename}{name}
\newcommand{\thesiscandidatesurname}{surname}
\newcommand{\thesissupervisoronetitle}{prof.}
\newcommand{\thesissupervisoronename}{name}
\newcommand{\thesissupervisoronesurname}{surname}
\newcommand{\thesissupervisortwotitle}{prof.}
\newcommand{\thesissupervisortwoname}{name}
\newcommand{\thesissupervisortwosurname}{surname}
\newcommand{\thesissupervisorthreetitle}{prof.}
\newcommand{\thesissupervisorthreename}{name}
\newcommand{\thesissupervisorthreesurname}{surname}
\newcommand{\thesissupervisorfourtitle}{prof.}
\newcommand{\thesissupervisorfourname}{name}
\newcommand{\thesissupervisorfoursurname}{surname}
\newcommand{\thesisdate}{DECEMBER 2024}
\newcommand{\thesiscourse}{COMPUTER ENGINEERING}
\newcommand{\thesisuniversity}{POLITECNICO DI TORINO}
\newcommand{\thesislevel}{MASTER} % master or bachelor
\newcommand{\thesiscandidatetext}{Candidate}
\newcommand{\thesissupervisortext}{Supervisors}

% fontsize is {size}{spacing}\family
\newcommand {\institutionfont}{\fontsize {22}{30}\scshape}
\newcommand {\divisionfont}{\fontsize {16}{20}\rmfamily}
\newcommand {\pretitlefont}{\fontsize {16}{16}\rmfamily}
\newcommand {\customtitlefont}{\fontsize {15}{18}\scshape}% {iwona}{bx}{n}}
\newcommand {\fixednamesfont}{\fontsize {12}{18}\mdseries}
\newcommand {\namesfont}{\fontsize {6.5}{11}\bfseries}
\newcommand {\candnamesfont}{\fontsize {10}{11}\bfseries}
\newcommand {\footfont}{\fontsize {15}{18}\rmfamily}

\setstretch{1.2}
%\onehalfspacing
%\doublespacing


\addbibresource{bibliography.bib}

% to load the glossaries, not needed if using bib2gls
% for glossary entry
% @entry{bird,
%     name={bird},
%     description = {feathered animal},
%     see={[see also]{duck,goose}}
% }

% if this bib file does not work, try using \input{file.tex}
% where all the \newabbreviation commands have been inserted
% containing all the definitions

% Gls to capitalize first letter
% GLS for full uppercase
% for abbreviations also
% glsxtrshort for abbreviation
% similar for long, full, and capital configurations, add pl at the end for plurals
% glsentryshort, long, plural (referred to shorts) must be used when in section titles
% glslink to allow the link but use a different text (as for href)


% if you want to use also description for the abbreviations/acronyms, you should use bib2gls and define all the entries in a bib file, which is incompatible with Overleaf
\newacronym{USB}{USB}{Universal Serial Bus}
\makeglossaries

\pagestyle{fancy}
\fancyhf{}  % Svuota le intestazioni esistenti
\fancyhead[C]{\itshape \leftmark}  % Titolo del capitolo in corsivo
\fancyfoot[C]{\thepage}  % Numero di pagina centrato in basso
% Personalizza \chaptermark per mostrare solo il titolo, senza "Chapter X"
\renewcommand{\chaptermark}[1]{\markboth{#1}{}}

\renewcommand{\thepage}{\Roman{page}}

\begin{document}

\title{\vspace*{-5mm}\textbf{\thesistitle}\\Summary} % vspace is needed to shift upwards the title
\date{\thesisdate}
\author{\textbf{Candidate}:\\\thesiscandidatename~\thesiscandidatesurname\\
\textbf{Supervisors}:\\\thesissupervisoronetitle~\thesissupervisoronename~\thesissupervisoronesurname\\
\thesissupervisortwotitle~\thesissupervisortwoname~\thesissupervisortwosurname\\
\thesissupervisorthreetitle~\thesissupervisorthreename~\thesissupervisorthreesurname}


%\ateneo{\thesisuniversity} % university name
\logosede[5cm]{\thesisuniversitylogo} % logo, square brackets contain the height

\titolo{\thesistitle} % title
%\sottotitolo{Metodo dei satelliti medicei} % subtitle

% place/remove a slash \\ to put the name on the following line or after Master Degree Course
\corsodilaurea{\thesiscourse} % course name


%~251197 % id number is not needed

\candidato{\thesiscandidatename~\textsc{\thesiscandidatesurname}} % candidate

% using tabular we can have more than 1 supervisor under the same column
\relatore{\tabular{@{}l}%
    \xmakefirstuc{\thesissupervisoronetitle}~\thesissupervisoronename~\textsc{\thesissupervisoronesurname}\\[0.4ex]
    \xmakefirstuc{\thesissupervisortwotitle}~\thesissupervisortwoname~\textsc{\thesissupervisortwosurname}\\[0.4ex]
    \xmakefirstuc{\thesissupervisorthreetitle}~\thesissupervisorthreename~\textsc{\thesissupervisorthreesurname}
    \endtabular}
%\terzorelatore{Ciao}

% in this way we have Academic Year without stile=classica, so without lines
%\sedutadilaurea{\textsc{Academic~Year} 2019-2020}% per la laurea magistrale
% for PoliTo there is only month year
\sedutadilaurea{\thesisdate}% per la laurea magistrale


\english

\iflanguage{english}{%
	%\retrofrontespizio{This work is subject to the Creative Commons Licence}

	\CorsoDiLaureaIn{\thesislevel's Degree Course in\space}

	\TesiDiLaurea{\thesislevel's Degree Thesis}

	\InName{in}
	\CandidateName{\xmakefirstuc{\thesiscandidatetext}}% or Candidates
	\AdvisorName{\xmakefirstuc{\thesissupervisortext}}% or Supervisor
	%\TutorName{Tutor}
	%\NomeTutoreAziendale{Internship Tutor}

	%\NomePrimoTomo{First volume}
	%\NomeSecondoTomo{Second Volume}
	%\NomeTerzoTomo{Third Volume}
	%\NomeQuartoTomo{Fourth Volume}
}{}

% front page
% frontespizio can be used for the first page print
% while the custom-made frontpage can be used as hard-cover
% use pdfjoin or pdfseparate to extract or put together the pages if needed
%\frontespizio* % without star the logo is on top
\newgeometry{top=4cm,left=3cm,right=3cm,bottom=4cm,heightrounded}
\begin{titlepage}
\centering
%
{\institutionfont \textbf{\MakeUppercase{\thesisuniversity}} \par}
%
\vspace{\stretch{2}} % changing this number and the others changes the proportion
%
{\divisionfont \textbf{\thesislevel's Degree in \thesiscourse} \par}
%
\vspace{\stretch{3}}
%

\includegraphics[width=50mm]{\thesisuniversitylogo}\\
%
\vspace{\stretch{4}}
%
{\divisionfont \textbf{\thesislevel's Degree Thesis} \par}
%
\vspace{\stretch{3}}
%
{\customtitlefont \textbf{\thesistitle} \par}
%
\vspace{\stretch{5}}
%
\makebox[\textwidth]{\null\hfill\def\arraystretch{2}% % to change the spacing change this number
\begin{minipage}[t]{.35\textwidth}\raggedright
    %\begin{adjustbox}{width={\textwidth},totalheight={\textheight},keepaspectratio} % with adjustbox it adapts to the lengths of the names, remove it if you want the same font dimension
    \begin{tabular}[t]{@{}l@{}}
        \fixednamesfont \textbf{\thesissupervisortext} \\
        \namesfont \xmakefirstuc{\thesissupervisoronetitle}~\thesissupervisoronename~\MakeUppercase{\thesissupervisoronesurname}\\
        \namesfont \xmakefirstuc{\thesissupervisortwotitle}~\thesissupervisortwoname~\MakeUppercase{\thesissupervisortwosurname}\\
        \namesfont \xmakefirstuc{\thesissupervisorthreetitle}~\thesissupervisorthreename~\MakeUppercase{\thesissupervisorthreesurname}\\
        \namesfont\xmakefirstuc{\thesissupervisorfourtitle}~\thesissupervisorfourname~\MakeUppercase{\thesissupervisorfoursurname}
    \end{tabular}
    %\end{adjustbox}
\end{minipage}
%
\hfill
%
\begin{minipage}[t]{.35\textwidth}\raggedleft
%\begin{adjustbox}{width={\textwidth},totalheight={\textheight},keepaspectratio} % with adjustbox it adapts to the lengths of the names, remove it if you want the same font dimension
\begin{tabular}[t]{@{}l@{}}
    \fixednamesfont \textbf{\thesiscandidatetext} \\
    \candnamesfont \thesiscandidatename~\MakeUppercase{\thesiscandidatesurname}
\end{tabular}
%\end{adjustbox}
\end{minipage}\hfill\null}\\
%
\vspace{\stretch{5}}
%
{\footfont \textbf{\thesisdate} \par}
%
\end{titlepage}

\restoregeometry
 % custom frontpage
%\retrofrontespizio
% insert text for the back of the front page
% if you insert any remove the following \paginavuota
% either a blank page or a back is needed to have double-sided printing
% pay attention to leave the space for the page

\paginavuota % clears a page

%\frontmatter

% abstract if needed
% \begin{abstract}
%     

%\fancyhf{} % clear all header and footer fields
%\fancyhead[RO,R]{\thepage} %RO=right odd, RE=right even
%\renewcommand{\headrulewidth}{0pt}

\begin{center}
    %\Large
    \textbf{Thesis Title}
    
    \vspace{0.4cm}
    
    \vspace{0.4cm}
    \textbf{Thesis Author}
    
    \vspace{0.9cm}
    \textbf{Abstract}
    
\end{center}

Lorem ipsum dolor sit amet, consectetur adipiscing elit, sed do eiusmod tempor incididunt ut labore et dolore magna aliqua. Ut enim ad minim veniam, quis nostrud exercitation ullamco laboris nisi ut aliquip ex ea commodo consequat. Duis aute irure dolor in reprehenderit in voluptate velit esse cillum dolore eu fugiat nulla pariatur. Excepteur sint occaecat cupidatat non proident, sunt in culpa qui officia deserunt mollit anim id est laborum.
% \end{abstract}

% to create blank pages for openright in frontmatter
% use one of the following two methods
% 1) use the following three lines
%\phantom{0} % needed otherwise cleardoublepage does not clean the page because it sees it empty
%\cleardoublepage
%\thispagestyle{empty} % to have empty page, without numbers
% 2) or
\paginavuota % to manually create a blank page

\pagenumbering{Roman}
%\sommario


%\fancyhf{} % clear all header and footer fields
%\fancyhead[RO,R]{\thepage} %RO=right odd, RE=right even
%\renewcommand{\headrulewidth}{0pt}

\begin{center}
    %\Large
    \textbf{Thesis Title}
    
    \vspace{0.4cm}
    
    \vspace{0.4cm}
    \textbf{Thesis Author}
    
    \vspace{0.9cm}
    \textbf{Abstract}
    
\end{center}

Lorem ipsum dolor sit amet, consectetur adipiscing elit, sed do eiusmod tempor incididunt ut labore et dolore magna aliqua. Ut enim ad minim veniam, quis nostrud exercitation ullamco laboris nisi ut aliquip ex ea commodo consequat. Duis aute irure dolor in reprehenderit in voluptate velit esse cillum dolore eu fugiat nulla pariatur. Excepteur sint occaecat cupidatat non proident, sunt in culpa qui officia deserunt mollit anim id est laborum.

\phantom{0}
\cleardoublepage
\thispagestyle{empty}

%\ringraziamenti
% acknowledgements

ACKNOWLEDGMENTS

\vspace*{5\baselineskip}

\begin{flushright}
    \textit{to...}
\end{flushright}


\tableofcontents


\listoffigures


\listoftables

% actually abbreviation is the name used for acronym in glossaries-extra
% title sets the name
% type tells the type of glossary to print
% style overrides the global style
% here we are printing only abbreviations
% printunsrtglossary if using record, otherwise printglossary is ok
\paginavuota
\printunsrtglossary[style=long,title=Acronyms,type=\acronymtype]

% also list of symbols here if needed

% to remove all first use occurrences given the presence of the summary
\glsresetall
% to skip all the first use occurrences, using only short forms
% \glsunsetall


%\mainmatter

%\part{Prima Parte} % parts division, not needed
% Chapters always open on a right-side page, i.e. odd numbers, so a blank page is inserted if needed
%\cleardoublepage % to have a fully blank page
% a blank page appears before the first chapter in some configurations, on the last version it doesn't

\paginavuota
\pagenumbering{arabic}

% list here all the chapters
\chapter{Chapter One}

\section{Section name}
Lorem ipsum dolor sit amet, consectetur adipiscing elit, sed do eiusmod tempor incididunt ut labore et dolore magna aliqua. Ut enim ad minim veniam, quis nostrud exercitation ullamco laboris nisi ut aliquip ex ea commodo consequat. Duis aute irure dolor in reprehenderit in voluptate velit esse cillum dolore eu fugiat nulla pariatur. Excepteur sint occaecat cupidatat non proident, sunt in culpa qui officia deserunt mollit anim id est laborum. \cite{IEEEexample:article_typical}

\begin{figure}
    \centering
    \includegraphics[width=0.5\linewidth]{00_drafting/images/polito_logo_2021_blu_cut.jpg}
    \caption{Image Caption}
    \label{fig:enter-label}
\end{figure}


\begin{longtable}{ | m{4em} | m{4em} | m{4em} | m{4em} | m{4em} | m{4em} | m{4em} | }
\hline
\textbf{Col name} & \textbf{Col name} & \textbf{Col name} & \textbf{Col name} & \textbf{Col name} & \textbf{Col name} & \textbf{Col name} \\
\hline
\endfirsthead
\hline
\textbf{Col name} & \textbf{Col name} & \textbf{Col name} & \textbf{Col name} & \textbf{Col name} & \textbf{Col name} & \textbf{Col name} \\
\hline
\endhead

text & text & text & text & text & text & text \\
\hline


\caption{Long Table}
\label{tab:GPS_ICD} \\
\end{longtable}

\chapter{Chapter Two}

\section{Section name}
Lorem ipsum dolor sit amet, consectetur adipiscing elit, sed do eiusmod tempor incididunt ut labore et dolore magna aliqua. Ut enim ad minim veniam, quis nostrud exercitation ullamco laboris nisi ut aliquip ex ea commodo consequat. Duis aute irure dolor in reprehenderit in voluptate velit esse cillum dolore eu fugiat nulla pariatur. Excepteur sint occaecat cupidatat non proident, sunt in culpa qui officia deserunt mollit anim id est laborum.

% \paginavuota % it works even without stile=classica

\appendix
\chapter{Appendix A}\label{ch:apda}

\markboth{Appendix A}{}

Lorem ipsum dolor sit amet, consectetur adipiscing elit, sed do eiusmod tempor incididunt ut labore et dolore magna aliqua. Ut enim ad minim veniam, quis nostrud exercitation ullamco laboris nisi ut aliquip ex ea commodo consequat. Duis aute irure dolor in reprehenderit in voluptate velit esse cillum dolore eu fugiat nulla pariatur. Excepteur sint occaecat cupidatat non proident, sunt in culpa qui officia deserunt mollit anim id est laborum.
\chapter{Appendix B}\label{ch:apda}

\markboth{Appendix B}{}

Lorem ipsum dolor sit amet, consectetur adipiscing elit, sed do eiusmod tempor incididunt ut labore et dolore magna aliqua. Ut enim ad minim veniam, quis nostrud exercitation ullamco laboris nisi ut aliquip ex ea commodo consequat. Duis aute irure dolor in reprehenderit in voluptate velit esse cillum dolore eu fugiat nulla pariatur. Excepteur sint occaecat cupidatat non proident, sunt in culpa qui officia deserunt mollit anim id est laborum.

% endnotes here if needed

\phantom{0}
\cleardoublepage
\printbibliography[heading=bibintoc] % heading required to show it in ToC


\Dedications
\pagestyle{empty} 
\vspace*{4\baselineskip}

\vspace{\baselineskip}

Dedications for everyone

\end{document}
